\documentclass[11pt]{article}
\usepackage[T1]{fontenc}
\usepackage{lmodern}
\usepackage[margin=1in]{geometry}
\usepackage{amsmath,amssymb,amsthm,mathtools}
\usepackage{graphicx}
\usepackage[hidelinks]{hyperref}
\usepackage{enumitem}
\emergencystretch=1.5em

\newtheorem{theorem}{Theorem}
\newtheorem{corollary}[theorem]{Corollary}
\newtheorem{lemma}[theorem]{Lemma}
\newtheorem{remark}[theorem]{Remark}
\newcommand{\cO}{\mathcal O}
\newcommand{\cA}{\mathfrak A}
\newcommand{\cM}{\mathcal M}
\newcommand{\cL}{\mathcal L}

\title{A Non-Archimedean Operator Algebra\
with Type-II Baer Reduction}
\author{Run \texttt{fa\_banach\_001}, lane 14}
\date{August 9, 2026}

\begin{document}
\maketitle

\begin{abstract}
Kochubei asked for a non-Archimedean operator algebra whose reduction is
a type-II Baer ring.  We give a general equal-characteristic realization
principle: every unital algebra $R$ over a field $k$ is the reduction of a
norm-closed unital operator algebra on a space $c_0(J,k((t)))$.  Choose a
Hamel basis of $R$ containing $1$, put the Gauss norm on
$k((t))\otimes_k R$, complete it, and use the left regular
representation.  The trivial norm on the coefficient field makes the
Gauss norm submultiplicative, while evaluation at $1$ makes the regular
representation isometric.  A $c_0$-family has only finitely many
coefficients with nonzero residue, so completion introduces no new
residue elements.  Applying the construction to von Neumann's algebraic
continuous factor $\cM_k$, a type-$\mathrm{II}_f$ regular self-injective
ring, produces the required type-II Baer reduction.
\end{abstract}

\section{Source problem and scope}

Let $K$ be a complete nontrivially valued non-Archimedean field, let
$\cO_K$ be its valuation ring, $\mathfrak m_K$ its maximal ideal, and
$\widetilde K=\cO_K/\mathfrak m_K$ its residue field.  For a norm-closed
unital algebra $\cA\subseteq\cL(c_0(J,K))$, write
\[
 \cA_1=\{T\in\cA:\|T\|\leq 1\},\qquad
 \cA_0=\{T\in\cA:\|T\|<1\},\qquad
 \widetilde{\cA}=\cA_1/\cA_0.
\]
Kochubei calls $\cA$ an operator algebra with Baer reduction when
$\widetilde{\cA}$ is a Baer ring \cite{Kochubei}.  At the end of
Section~4, on page~7, he states:
\begin{quote}
``The problem of constructing non-Archimedean operator algebras with
type II Baer reductions remains open.''
\end{quote}

\begin{figure}[ht]
\centering
\includegraphics[width=\textwidth]{figures/open_problem_crop.png}
\caption{The open-problem statement and its immediate context in
Kochubei, page~7 \cite{Kochubei}.}
\end{figure}

The result below answers the literal existence problem in the general
class of complete non-Archimedean fields used in the paper's definition.
It gives an equal-characteristic example and does not assert that the
example is a group algebra.  Section~4 of the source paper begins with the
temporary convention $K=\mathbb C_p$.  If the final sentence is instead
read as a problem specifically over that fixed mixed-characteristic
field, our construction is a substantial partial answer but does not
settle that narrower interpretation; see Section~\ref{sec:scope}.

\section{Proof intuition}

The key point is that equal characteristic provides a coefficient field.
Take $K=k((t))$, so that $k$ embeds in the valuation ring and maps
isomorphically onto the residue field.  If $E$ is a Hamel basis of an
arbitrary $k$-algebra $R$, every product of two basis elements is a finite
$k$-linear combination of basis elements.  Since every nonzero scalar in
$k$ has $t$-adic norm one, the resulting structure constants are uniformly
bounded by one.  Thus the ordinary Gauss norm on $K\otimes_kR$ is
submultiplicative and its completion is the $c_0$-space on $E$.

Left multiplication acts contractively on this completion.  Because the
basis contains $1$, applying left multiplication by $a$ to $1$ recovers
$a$ and proves that its operator norm is exactly the Gauss norm.  The
regular representation is therefore a closed operator algebra.  Finally,
an element of the unit ball of $c_0(E,K)$ has only finitely many
coefficients of norm one; reducing those coefficients gives an element of
the original algebra $R$.  The kernel consists exactly of the elements of
norm strictly less than one.  Consequently the operator-algebra reduction
is $R$, with no enlargement caused by completion.

\section{An equal-characteristic realization theorem}

\begin{theorem}[Realization theorem]\label{thm:realization}
Let $k$ be any field and let $R$ be any unital associative $k$-algebra.
There exist
\[
 K=k((t)),\qquad B=c_0(E,K),\qquad
 \cA\subseteq\cL(B),
\]
where $K$ has its $t$-adic absolute value, $E$ is a Hamel basis of $R$ over
$k$, and $\cA$ is a norm-closed unital $K$-operator algebra, such that
\[
 \widetilde{\cA}\cong R
\]
as unital $k$-algebras.
\end{theorem}

\begin{proof}
Choose a Hamel basis $E$ of $R$ containing $1_R$.  On the algebraic scalar
extension
\[
 A_{\mathrm{alg}}=K\otimes_kR
 =\left\{\sum_{e\in F}a_e e:F\subset E\text{ finite},\ a_e\in K\right\}
\]
put the Gauss norm
\[
 \left\|\sum_{e\in F}a_e e\right\|_G=\max_{e\in F}|a_e|.
\]
For $e,f\in E$, write
\[
 ef=\sum_{g\in F(e,f)}c_{e,f}^g g,
 \qquad c_{e,f}^g\in k,
\]
with $F(e,f)$ finite.  The $t$-adic absolute value is trivial on $k$, so
$|c_{e,f}^g|\leq1$.  The ultrametric inequality therefore gives
\begin{equation}\label{eq:submult-alg}
 \|xy\|_G\leq\|x\|_G\|y\|_G
 \qquad(x,y\in A_{\mathrm{alg}}).
\end{equation}

The completion $A$ of $A_{\mathrm{alg}}$ is, as a Banach $K$-space,
canonically $c_0(E,K)$.  Inequality \eqref{eq:submult-alg} extends the
multiplication uniquely to $A$ and gives
\begin{equation}\label{eq:submult-complete}
 \|xy\|\leq\|x\|\|y\|\qquad(x,y\in A).
\end{equation}
Thus $A$ is a unital non-Archimedean Banach algebra and $\|1_R\|=1$.

For $a\in A$, let $L_a\in\cL(A)$ be left multiplication by $a$.
Equation \eqref{eq:submult-complete} gives $\|L_a\|\leq\|a\|$.  The
reverse inequality follows by evaluating at the unit:
\[
 \|L_a\|\geq\|L_a(1_R)\|=\|a\|.
\]
Hence $a\mapsto L_a$ is an isometric unital algebra homomorphism.  Its
image
\[
 \cA=\{L_a:a\in A\}\subseteq\cL(c_0(E,K))
\]
is consequently norm closed.

It remains to identify the reduction.  If
$a=\sum_{e\in E}a_e e\in A$ and $\|a\|\leq1$, then $(a_e)_{e\in E}$ is a
$c_0$-family in $\cO_K$.  In particular, only finitely many coefficients
satisfy $|a_e|=1$, so
\begin{equation}\label{eq:rho}
 \rho(a)=\sum_{e\in E}\overline{a_e}\,e
\end{equation}
is a well-defined finite sum in $R$.  Reduction of the structure
constants shows that $\rho:A_1\to R$ is a surjective unital algebra
homomorphism.

We claim that $\ker\rho=A_0$.  One inclusion is immediate.  Conversely,
if $\rho(a)=0$, then $|a_e|<1$ for every $e$.  Since $(a_e)$ is a
$c_0$-family, the set
\[
 F=\{e\in E:|a_e|\geq 1/2\}
\]
is finite.  If $F$ is nonempty, the maximum of $|a_e|$ over $F$ is
strictly below one, while outside $F$ all coefficients have norm below
one half.  Hence $\|a\|<1$.  If $F$ is empty, directly
$\|a\|\leq1/2<1$.  This proves the claim and yields $A_1/A_0\cong R$.

The regular representation is isometric, so the unit balls and open unit
balls of $A$ correspond exactly to those of $\cA$.  Hence
$\widetilde{\cA}\cong A_1/A_0\cong R$, as required.
\end{proof}

\section{The type-II example}

For a field $k$, let $\cM_k$ denote von Neumann's algebraic continuous
factor: the completion, in its unique rank metric, of
\[
 \varinjlim_n M_{2^n}(k),\qquad x\longmapsto
 \begin{pmatrix}x&0\\0&x\end{pmatrix}.
\]
Ara and Claramunt recall that $\cM_k$ is a simple regular right-and-left
self-injective ring and a continuous factor of type $\mathrm{II}_f$
\cite[Introduction and Section~2]{AraClaramunt}.  A regular self-injective
ring is Baer, and the type $\mathrm{II}_f$ classification says precisely
that its unit is finite and it has no nonzero Abelian idempotent; this is
type II in Kochubei's terminology \cite{Goodearl,Kochubei}.

\begin{corollary}[Type-II Baer reduction]\label{cor:typeII}
For every field $k$ there is a complete nontrivially valued
non-Archimedean field $K=k((t))$, a Banach space $B$ over $K$ with an
orthonormal basis, and a norm-closed unital operator algebra
$\cA\subseteq\cL(B)$ such that $\widetilde{\cA}$ is a type-II Baer ring.
In fact one may arrange
\[
 B=c_0(E,K),\qquad \widetilde{\cA}\cong\cM_k.
\]
\end{corollary}

\begin{proof}
Apply Theorem~\ref{thm:realization} to $R=\cM_k$.
\end{proof}

This proves the requested existence statement.  Notice that the
construction is stronger at the realization level: the analytic step does
not use any special feature of a continuous factor and can realize every
unital $k$-algebra as a reduction.

\section{Verification and limitations}\label{sec:scope}

The proof has four independently checkable points.
\begin{enumerate}[leftmargin=*]
\item Every basis product has finite support, and its coefficients lie in
the trivially valued subfield $k$; this proves the Gauss-norm estimate.
\item Completion is $c_0(E,K)$, and the bounded multiplication extends
uniquely.
\item The left regular representation is isometric because $1_R$ is a
basis vector of norm one; hence the image is a closed operator algebra.
\item A $c_0$-family has finite residue support, and the kernel argument
after \eqref{eq:rho} proves that no extra residue elements appear.
\end{enumerate}
No computational experiment is used as evidence.  The only external
structural input is the standard existence and type-$\mathrm{II}_f$
classification of $\cM_k$.

The construction does not produce a group algebra.  It also uses a
coefficient field $k\hookrightarrow k((t))$ and therefore does not directly
give an example over $\mathbb C_p$.  The source paper's general definition
allows arbitrary complete nontrivially valued $K$, while its final examples
are presented under the Section~4 convention $K=\mathbb C_p$.  Thus the
packet's status is: candidate full solution for the literal
field-unrestricted existence problem; candidate partial result if a
reviewer reads the sentence as imposing the standing field
$K=\mathbb C_p$.

\paragraph{Bounded novelty search.}
On August 9, 2026, the run's registry, solution, attempt, and proof-gap
indexes were searched for arXiv:1401.2268 and the phrases ``type II Baer
reduction'', ``Baer reductions'', and ``non-Archimedean operator algebra''.
The local arXiv full-text corpus was searched for the same phrases.  Web
searches used the exact source title, the exact open-problem sentence,
Kochubei's name, and the same core phrases.  These searches found the
source paper and background work but no later paper claiming a type-II
construction or an answer to this problem.  Ara--Claramunt
\cite{AraClaramunt} supplies the algebraic continuous factor used here but
does not identify the non-Archimedean realization or Kochubei's question.
Accordingly, novelty is plausible but not certified by an exhaustive
bibliographic database search.

\paragraph{Human-review recommendation.}
The candidate should receive expert review.  The main proof is short; the
highest-value checks are the equality $\ker\rho=A_0$, compatibility of the
two type-II conventions, and whether the source's final sentence is meant
field-unrestricted or specifically over $\mathbb C_p$.

\begin{thebibliography}{9}
\bibitem{AraClaramunt}
P.~Ara and J.~Claramunt,
\emph{Uniqueness of the von Neumann continuous factor},
Canad. J. Math. \textbf{70} (2018), 961--982;
\href{https://arxiv.org/abs/1705.04501}{arXiv:1705.04501}.

\bibitem{Goodearl}
K.~R. Goodearl,
\emph{Von Neumann Regular Rings},
Pitman, London, 1979.

\bibitem{Kochubei}
A.~N. Kochubei,
\emph{Non-Archimedean group algebras with Baer reductions},
Algebr. Represent. Theory \textbf{17} (2014), 1861--1867;
\href{https://arxiv.org/abs/1401.2268}{arXiv:1401.2268}.
\end{thebibliography}

\end{document}
